\title{LongRoPE: Extending LLM Context Window Beyond 2 Million Tokens}

\begin{abstract}
An extensive context window is highly sought after in large language models (LLMs). Nonetheless, existing models are constrained to a maximum of approximately 128k tokens, primarily because of prohibitive fine-tuning expenses, limited access to sufficiently long training samples, and disruptive outlier values generated by previously unseen token positions.
\end{abstract}

This study presents LongRoPE, marking the first instance of extending the context window in pre-trained LLMs to a remarkable \textbf{2048k} tokens. This expansion is accomplished using as few as 1k fine-tuning steps restricted to a 256k token context length during training, all while preserving model efficacy for short context windows. LongRoPE is underpinned by three main contributions: (i) we discover and leverage two distinct non-uniformities in positional interpolation through a systematic and efficient search, yielding a superior initialization point for fine-tuning, and enabling an 8$\times$ expansion even without fine-tuning; (ii) we implement a staged extension method, first fine-tuning the LLM up to 256k tokens, followed by a subsequent positional interpolation on this adapted model to realize a 2048k token context window; (iii) we recalibrate LongRoPE at 8k tokens to restore performance for short context tasks. Through comprehensive experiments conducted on LLaMA2 and Mistral across multiple benchmarks, we validate the efficacy of the proposed approach. Crucially, models extended with LongRoPE maintain the original network design, only introducing minimal modifications to the positional embeddings, and retain compatibility with most existing system optimizations. Code resources will be made available at \texttt{https://github.com/microsoft/LongRoPE}

\section{Introduction}

Although Large Language Models (LLMs) have achieved impressive results across a range of tasks~\cite{gpt4,llama2}, they are frequently constrained by a limited context window—such as the 4096 token restriction in LLaMA2~\cite{llama2}. When required to process input exceeding this limit, the models exhibit a drop in performance, which is attributed to encountering positional tokens outside their training distribution. This limitation creates difficulties in critical applications, including in-context learning with large numbers of examples~\cite{Huang2023FewerIM} and the effective deployment of LLM agents~\cite{park2023generative,madaan2023selfrefine}.

Recent research has demonstrated that the context window of a pre-trained LLM can be increased to approximately 128k tokens by fine-tuning with longer documents~\cite{longlora,pi,yarn,zhang2024soaring,liu2023scaling}. However, pushing these limits presents three significant challenges. First, adding previously untrained position indices introduces a large number of anomalous values, resulting in out-of-distribution complications and hindering convergence during fine-tuning~\cite{pi}. This complication is amplified when the context window is expanded from 4k to over 1000k, as such an extension brings in more than 90\% unfamiliar positions. Second, the fine-tuning process generally relies on the availability of input texts that match the intended window size, yet there is a shortage of such long documents, particularly those surpassing 1000k in length, in existing data collections. Additionally, processing these ultra-long texts during training demands substantial computational power, with enormous training time and GPU consumption. Third, as the context window is further extended, the model’s attention mechanism is forced to allocate its capacity over a vastly larger sequence, causing attention to be thinly spread over numerous positions and resulting in decreased effectiveness for tasks with shorter context lengths~\cite{pi}.

A common strategy for addressing the initial challenge is to interpolate RoPE positional embeddings~\cite{rope,pi}, a technique that rescales incoming positional indices to fall within the pretrained interval, as illustrated in Fig.\ref{fig:overview}. Position Interpolation (PI)~\cite{pi} achieves this by linearly scaling the rotational angles of RoPE according to the proportion of extension. In contrast, NTK~\cite{ntk,dynamicntk} proposes using uneven interpolation and extrapolation across the dimensions of RoPE. YaRN~\cite{yarn} takes a different route by separating RoPE's dimensions into three groups based on frequency, then employing extrapolation, NTK, and linear interpolation methods to each group individually. Despite these advancements, the positional embedding layer in Transformer models contains intricately non-uniform information entropy. Current methodologies largely overlook this subtle heterogeneity, resulting in information degradation and thereby restricting the maximum achievable context window.

Section~\ref{sec:analysis} provides empirical evidence for two principal insights: \textbf{(1)} Successful positional interpolation must address two kinds of non-uniformity: fluctuations in RoPE dimensions and token positions. Reduced RoPE dimensionality and early token positions require minimal interpolation, although the best strategy depends on the intended extrapolation length. \textbf{(2)} Incorporating these non-uniformities into the positional interpolation process helps preserve the integrity of the original RoPE representation, especially across critical dimensions and token positions. This approach reduces information loss during interpolation, thus yielding improved initial weights for fine-tuning. Additionally, in settings where fine-tuning is not used, this method facilitates up to an 8$\times$ extension.

Building on these observations, we propose \textbf{LongRoPE}, a robust approach designed to increase the LLM context window to more than 2 \emph{million} tokens. The LongRoPE framework incorporates three principal contributions. First, {LongRoPE} takes comprehensive advantage of multidimensional variations in positional interpolation by determining optimal rescale factors for the rotational angles in RoPE, tailored for each RoPE dimension and grounded in the specific token positions. Because the parameter space for identifying these rescale factors increases exponentially with the desired context window expansion, {LongRoPE} utilizes an evolutionary search algorithm complemented by two search optimization methods to enhance efficiency. Figure~\ref{fig:overview} illustrates a representative instance of the optimized rescaled RoPE.

Subsequently, {LongRoPE} adopts a resource-efficient, staged extension method to reach a 2048k context window, circumventing the necessity for direct fine-tuning on exceedingly lengthy text corpora, which are both rare and difficult to obtain. This approach initiates with a search at the 256k sequence length using the pre-trained LLM, followed by fine-tuning the model at this sequence length. Leveraging the non-uniform positional interpolation mechanism—enabling up to an 8$\times$ expansion without additional fine-tuning—we perform a subsequent search to determine appropriate RoPE rescale factors on the previously fine-tuned, extended LLM. This process ultimately allows the realization of the 2048k context window for both LLaMA2 and Mistral~\cite{mistral}.

In order to address potential declines in performance when using the original (shorter) context lengths, {LongRoPE} also revises the RoPE rescale factors for the LLM that has been extended. Analogous to expanding the context window from 256k to 2048k, we employ our search algorithm to scale back the window to 4k and 8k on the LLM previously fine-tuned at 256k. This adjustment is intended to reduce the degree of positional interpolation. When the input sequence during inference is shorter than 8k, RoPE is updated with the rescale factors obtained through this search procedure.

A comprehensive set of experiments involving multiple LLM architectures and a variety of long-context evaluations highlights the efficacy of our approach. Our findings indicate that LongRoPE consistently sustains low perplexity across evaluation lengths ranging from 4k up to 2048k, secures passkey retrieval accuracy exceeding 90\%, and maintains accuracy on established benchmarks constructed within a 4096 token context window. The LongRoPE method is compatible with all LLMs utilizing RoPE embedding. We intend to make both our codebase and the LongRoPE-2048k models publicly available.

\section{Non-uniformity in Positional Interpolation}
\label{sec:analysis}

\subsection{Preliminary}

Transformer architectures necessitate the incorporation of explicit positional cues—typically implemented through position embedding—to encode the sequential order of input tokens. In this study, we direct our attention toward the RoPE~\cite{rope} position embedding technique, which has become a standard approach in current large language models (LLMs). Given a token located at position $n$, the RoPE encoding associated with it may be succinctly expressed as:

\begin{equation}
		\fontsize{7.8}{7.8} \selectfont
	\label{eq:rope}
	\left[ cos(n\theta_0), sin(n\theta_0), cos(n\theta_1), \cdot\cdot\cdot, cos(n\theta_{d/2-1}), sin(n\theta_{d/2-1}) \right]   
\end{equation}

Here, $d$ refers to the embedding dimension, 
 $n\theta_i$ denotes the rotary angle assigned to the token at sequence position $n$, and 
$\theta_i = \theta^{-2i/d}$ captures the frequency components for rotation. For RoPE, the standard value of $\theta$ is set to 10000.

\textbf{\emph{Context window extension ratio $s$} and positional interpolation}. We define $s$ as  the ratio of extended context length $L'$ to the original length $L$: $s=\frac{L'}{L}$. 

To expand the context window from $L$ to $L'$, prevailing positional interpolation techniques advocate reducing the rotation frequencies $\theta_i$ according to the extension ratio $s$. Defining $\beta=\theta^{2/d}$, and using $\lambda$ to represent the true scaling factor as a function of $s$, these various positional interpolation schemes can be expressed in a unified manner:

\begin{equation}	
	\fontsize{7.5}{7.5} \selectfont
	\label{eq:unified_rope}
	\begin{aligned}
		\left[cos\left(\frac{n}{\lambda(\beta)^0}\right), sin \left(\frac{n}{\lambda(\beta)^0}\right), cos\left(\frac{n}{\lambda(\beta)^1}\right), \cdot\cdot\cdot,  sin \left(\frac{n}{\lambda(\beta)^{d/2-1}}\right) \right]   
	\end{aligned}
\end{equation}

\textbf{Linear positional interpolation (PI)}. PI~\cite{pi} employs a strategy of linearly interpolating positional indices as long as they remain within the original pre-trained sequence length. To achieve an extension by a factor of $s$, the rotation angles corresponding to every position are scaled linearly by a factor of $\lambda = s$ for every RoPE dimension. This method, however, condenses position information, making it more challenging for the model to differentiate tokens that are close together in the sequence. Consequently, the performance of PI typically degrades when applied at larger extension ratios.

\textbf{NTK-based interpolation and extrapolation}. ~\cite{ntk,dynamicntk} examine RoPE through the lens of information encoding and utilize Neural Tangent Kernel (NTK) theory~\cite{jacot2018neural,tancik2020fourier}. To address the challenge of positional congestion in PI, their approach redistributes interpolation demands among RoPE’s various dimensions. In practice, lower (corresponding to high-frequency) dimensions experience less scaling, whereas upper (low-frequency) dimensions are scaled more aggressively. This technique enables both positional interpolation and extrapolation, characterized by $\lambda=s^{i}$. The enhanced dynamic NTK method~\cite{dynamicntk} further adapts the scaling factor at individual positions according to the specific sequence length encountered. In contrast to PI, which requires model fine-tuning for extension, NTK-aware approaches are capable of enlarging the context window even when fine-tuning is not performed, though this typically achieves at most a 4$\times$ extension ratio.

\textbf{YaRN}~\cite{yarn} presents a notable enhancement in positional interpolation by segmenting RoPE dimensions into three distinct groups according to their frequencies, and applying specific interpolation methods to each. In this approach, the high frequency dimensions are subject to extrapolation ($\lambda$=1), the low frequency components leverage linear interpolation (PI), and the intermediate frequency dimensions are managed via the NTK method. A central feature of YaRN is its method for grouping the RoPE dimensions; however, this classification is currently based on heuristic, manually-guided experimentation, which could potentially lead to less-than-optimal outcomes on newly developed LLMs.

\subsection{Study on Non-uniform Positional Interpolation}

Drawing inspiration from NTK and YaRN, we observe that their improvements are attributable to the introduction of non-linearity, particularly through the handling of various frequencies in RoPE dimensions to tailor interpolation and extrapolation. Nevertheless, present approaches to non-linearity predominantly depend on heuristics crafted by humans. This leads to two primary considerations: (1) Can the existing scheme for positional interpolation be further optimized? (2) Might there exist alternative non-linearities that remain unexplored?

To address these inquiries, we employ evolution search (refer to Sec.~\ref{sec:method}) to identify improved non-uniform positional interpolations tailored for LLaMA2-7B. Our search process leverages perplexity as the guiding metric, evaluating candidates using 5 randomly selected examples from the PG19~\cite{pg19} validation dataset. This empirical investigation leads us to uncover several important insights.

\textbf{Finding 1}: \textit{RoPE dimensions exhibit substantial non-uniformities, which are not effectively handled by current positional interpolation methods.}

We determine the optimal value of $\lambda$ for each RoPE dimension as described in Eq.~\ref{eq:unified_rope}. Table~\ref{tbl:exp1} presents a comparison of the perplexity scores achieved by LLaMA2-7B across various approaches on the PG19 and Proof-pile~\cite{proof-pile} test datasets, all evaluated in a zero-shot setting without additional fine-tuning. The results indicate that our approach yields clear improvements over existing methods. This suggests that both traditional linear interpolation (PI) and more recent non-uniform schemes (such as Dynamic-NTK and YaRN) do not provide optimal performance. Of particular note, YaRN performs worse than both PI and NTK on the PG19 dataset; this appears to be due to its failure to achieve the intended context window size when applied to large language models that have not undergone fine-tuning. For instance, YaRN's perplexity sharply increases beyond a context length of 7k in an 8k token setting.

Our investigation reveals that the rescaling parameters $\lambda$ in Eq.~\ref{eq:unified_rope} are heterogeneous, diverging from the uniform scaling factor $s$ used in PI, the constant value in NTK’s formula, or the block-based scaling applied by YaRN. Employing these variable rescaling factors yields considerable enhancements in the perplexity of LLaMA2 over extended context lengths of 8k and 16k tokens, all without the need for additional fine-tuning. The reason for this improvement is that the adapted positional embedding more faithfully maintains the properties of the original RoPE, particularly in the most influential dimensions, thereby lessening the model’s challenges when distinguishing between tokens in nearby positions.

\textbf{Finding 2}: \textit{RoPE for the initial tokens in the input sequence should be extrapolated with less interpolation.} 

We propose that for the first $\hat{n}$ tokens in the input sequence, it is preferable to minimize the use of RoPE interpolation. This suggestion arises because these tokens tend to attract higher attention scores and play a pivotal role within attention layers, as previously documented by Streaming LLM~\cite{streamingllm} and LM-Infinite~\cite{han2023lminfinite}. To empirically test this premise, we extend the context window to 8k and 16k utilizing PI and NTK, selectively withholding interpolation for the initial $\hat{n}$ tokens ($\hat n$ ranges over 0, 2, ..., 256). Setting $\hat{n}$=0 recovers the conventional behavior of PI and NTK. The results presented in Table~\ref{tbl:tokenposition} reveal two major findings: \textbf{(1)} Excluding position interpolation for the earliest tokens leads to enhanced performance. \textbf{(2)} The optimal value of $\hat{n}$, corresponding to the number of initial tokens left uninterpolated, varies with the desired context window extension.

\textbf{Finding 3}: \textit{Non-uniform positional interpolation effectively extends LLM context window in both fine-tuning and non-fine-tuning settings.} 

Although our investigation demonstrates that non-uniform position interpolation, as discovered through our search method, substantially enhances extension capabilities at 8k and 16k context windows even in the absence of fine-tuning, attaining robust performance at even greater lengths necessitates fine-tuning. Therefore, we conduct fine-tuning on LLaMA2-7B utilizing our optimized RoPE approach to accommodate a 64k context window (refer to the Appendix for detailed configurations). As presented in Table~\ref{tbl:finetune}, our approach yields significant improvements over both PI and YaRN, evident both prior to and following the fine-tuning process for LLaMA2-7B. This improvement can be attributed to the superior handling of non-uniform positional interpolation, which not only curtails information degradation but also establishes a stronger basis for subsequent fine-tuning.

\textbf{Summary}. This research highlights two core sources of non-uniformity: the distribution across RoPE dimensions and variability among token positions. Strategically leveraging these in positional interpolation leads to marked advancements in the ability of LLMs to manage extended context windows.

\section{LongRoPE}

\label{sec:method}
Building upon these insights, we propose LongRoPE. This approach initially develops an effective search strategy designed to maximize the utility of the identified non-uniformities, and subsequently leverages this technique to expand the context window of LLMs to over 2 million tokens.

\subsection{Problem Formulation}

The presence of these two non-uniformities results in an extensive solution landscape, complicating the optimization process. To manage this challenge, we recast the problem of optimizing multidimensional non-uniform position interpolation as a search-based problem.

Given an LLM with a context window size of $L'$ and input documents $\mathbf{X}$ containing sequences $\mathbf{x}\in\mathbf{X}$ that exceed $L'$ tokens, we represent the rotary angle associated with the $i^{th}$ dimension in RoPE embeddings at token index $n$ as $\frac{n}{\beta^i}$. Accordingly, the optimization objective is established as follows:

\begin{equation}
\begin{gathered}
		\label{eq:problem}

\underset {\mathbf{x}\in\mathbf{X}; \, |\mathbf{x}|\ge L'}{ \operatorname {arg\,min} } \,  \mathcal{L}\, \left( \text{LLM} (\text{RoPE}, \mathbf{X})\right), \text{where} 
\\
\scalebox{0.93}{
$\underset {\begin{subarray}{l} i=0,\cdot\cdot,\frac{d}{2}-1;\\ n\in[ 0,|\mathbf{x}|);\end{subarray}}{\text{RoPE}_(n)}= {\left[\cdots,\cos\left(\mathbb{I}(\hat{\lambda}_{i}, \hat{n})\cdot\frac{n}{\beta^i}\right), \sin\left(\mathbb{I}(\hat{\lambda}_{i}, \hat{n})\cdot\frac{n}{\beta^i}\right),\cdots\right] }$}\\
	\text{where}~\mathbb{I}(\hat{\lambda}_{i},\hat{n})= \begin{cases}
	1 & \text{if}~n<\hat{n} \\
	\frac{1}{\lambda_{i}} & \text{if}~n\geq\hat{n}
\end{cases}
	\end{gathered}
\end{equation}

In this context, we define a set of scaling factors, $\mathbb{I}(\hat{\lambda}_{i},\hat{n})$, to address the two types of non-uniformities present. Here, $\hat{\lambda_i}$ represents the dimension-specific non-uniformity within RoPE, and $\hat{n}$ refers to the non-uniformity across token positions. The scaling function $\mathbb{I}(\hat{\lambda}_{i},\hat{n})$ is applied to adjust the rotation angle associated with the $i^{th}$ RoPE dimension; the scaling factor for this dimension is given by $\hat\lambda_{i}$, and $\hat{n}$ acts as a position threshold. For tokens occupying the initial $\hat{n}-1$ positions, the scaling factor $\hat\lambda_{i}$ is inactive, such that the original RoPE rotation angle, $\frac{n}{\beta^i}$, is retained. Conversely, for token positions where $n \geq \hat{n}$, the scaling factor is enacted.

Assuming a desired context window of size $L'$, our task is to determine the most suitable set of rescale factors ($\mathbb{I}(\hat{\lambda}_{0},\hat{n})$, $\mathbb{I}(\hat{\lambda}_{1},\hat{n})$, ...$\mathbb{I}(\hat{\lambda}_{i},\hat{n})$...) across all RoPE dimensions from $1^{st}$ through $d^{th}$. The intention is for the modified $\text{LLM}$, leveraging the rescaled $\text{RoPE}$, to achieve the lowest possible next-token prediction loss $\mathcal{L}$ (namely, perplexity) on input sequences $\mathbf{X}$ of length $L'$.

\subsection{Searching the Non-uniform Position Interpolation}

To address the issue described in Eq.~\ref{eq:problem}, we present a straightforward yet powerful approach that seeks optimal rescaling coefficients for RoPE. This approach is designed to fully harness the multidimensional and non-uniform characteristics inherent in position embeddings.

\textbf{Search space}. We construct an expansive search space that explicitly incorporates both sources of non-uniformity. Table~\ref{tbl:searchspace} details the structure of this search space. More precisely, our method considers distinct rescale factors for every dimension within the RoPE embedding. For the sake of simplifying the construction of the search space, we focus on searching over $\lambda_i$ and $\hat{n}$, as opposed to directly searching for $\mathbb{I}(\hat{\lambda}_{i},\hat{n})$, where $\hat{\lambda}_{i}=1/\lambda_i$. According to Table~\ref{tbl:searchspace}, the parameter $\lambda_i$ can take values ranging from a lower bound of 1.0 (corresponding to straightforward extrapolation) up to an upper bound of $s\times1.25$ (indicating more extensive interpolation compared to PI), with increments of 0.01, where $s$ denotes the desired ratio for context window expansion.

The parameter $\hat{n}$ determines how many of the earliest token positions preserve the original RoPE embedding without position interpolation. In practice, we let $\hat{n}$ take values from the set \{0, 1, 2, 4, 8, 12, 16, 20, 24, 28, 32, 64, 128, 256\}. If $\hat{n}=0$, then every token position applies the identified rescale factors.

\textbf{Evolution-based search}. The search space outlined in Table~\ref{tbl:searchspace} encompasses a vast array of positional interpolation options, making thorough exploration computationally demanding. For instance, a $s=4\times$ extension yields $400^{128/2}\times14$ = $4\times10^{167}$ possible configurations. As the extension ratio increases, the complexity of the search space grows exponentially. To efficiently navigate this space, we employ evolution search~\cite{guo2020single} and propose two optimization strategies that markedly improve the efficiency of the search process. The complete search method is described in Algorithm~\ref{alg:search}.

\textit{Optimized initial population generation}. Rather than fully randomizing the initialization of a population of $P$ rescale factors, our method incorporates the three specific RoPE rescale factors associated with PI, NTK, and YaRN directly as initial members of the population. The rest of the individuals, amounting to $P-3$, are produced by applying random mutations—with a probability $p$—to these three core rescale factors.

\textit{Monotonically non-decreasing constraint}. Once the initial population is established, we evaluate each individual’s LLM perplexity. This involves applying the designated RoPE rescale factors to the target LLM and assessing the perplexity over the input $\mathbf{X}$. The $k$ individuals with the lowest perplexity scores are then selected as parents for subsequent evolutionary steps. However, due to the immense search space, unstructured mutation and crossover may result in exploration of suboptimal candidates, incurring significant costs in perplexity evaluation. This inefficiency is amplified for large $L'$, as each perplexity inference is computationally intensive.

To tackle this issue, we enforce a monotonic non-decreasing constraint on the sequence of sampled RoPE rescaling factors: $\lambda_i\le \lambda_{i+1}$. Only RoPE variants adhering to this constraint are utilized for perplexity assessment in LLMs, thereby substantially decreasing the search complexity. Concretely, this requirement ensures that $\lambda_i$ grows monotonically with respect to the RoPE dimension index ($i$ = 0,...,63). This form of dimensional monotonicity is inspired by NTK theory~\cite{jacot2018neural,tancik2020fourier,ntk}, which argues that lower-index (higher frequency) dimensions necessitate less interpolation (smaller $\lambda_i$), whereas higher-index (lower frequency) dimensions benefit from greater interpolation (larger $\lambda_i$).

\begin{algorithm}[t]
	\small
	\caption{The search algorithm}
	\textbf{Input:} target LLM, input samples $\mathbf{X}$,   population size $P$, mutation size $N_1$, crossover size $N_2$, max iterations $\mathcal{T}$, mutate probability $p$\\

	\begin{algorithmic}[1]
		\label{alg:search}
		\STATE Initialize Top-k as $\phi$;	
		\STATE $\text{P}_0 \leftarrow \textit{Initialize\_population\_with\_optimization}(P, \mathbf{X}, p)$; 
		\FOR{$i = 1$ to $\mathcal{T}$}
		\STATE \textit{Compute\_perplexity}(LLM, $\text{P}_{i-1}$, $\mathbf{X}$);
		\STATE  Top-k $\leftarrow$ \textit{Update\_Topk}(Top-k, $\text{P}_{i-1}$);
		\STATE $\text{P}_{mutation} \leftarrow \textit{Mutation\_with\_mono\_constraint}$(Top-k, $N_1$, $p$); 
		\STATE $\text{P}_{crossover} \leftarrow \textit{Crossover\_with\_mono\_constraint}$(Top-k, $N_2$);
		\STATE $\text{P}_i \leftarrow \text{P}_{mutation} \cup \text{P}_{crossover} \cup$ Top-k;
		\ENDFOR
		\STATE Output the individual from Top-k with the minimum perplexity;
	\end{algorithmic}
\end{algorithm}

\textbf{8$\times$ extension without fine-tuning}. Utilizing our evolutionary search approach, we successfully discover optimal non-uniform RoPE rescale factors that safeguard essential dimensions and token positions, thereby reducing information loss resulting from interpolation. As shown in Fig.\ref{fig:non-ft}, this technique allows for expanding the context window of LLaMA2 from 4k to 32k tokens with no additional fine-tuning. In comparison, prior methods like PI as well as non-uniform NTK and YaRN show a significant rise in perplexity when the context window is expanded beyond 2$\times$.

\subsection{Extending LLM Context Window to 2048K}
\label{sec:extendto2048k}

\textbf{Progressive extension to 2048k}. In this section, we present our approach for expanding the context window of pre-trained LLMs beyond the conventional 4k limit, reaching up to and beyond 2048k. Our results indicate that using non-uniform positional interpolation enables an 8$\times$ extension without the need for fine-tuning. However, when scaling up the context window by much larger factors (such as 512$\times$), fine-tuning becomes indispensable. One approach involves identifying appropriate RoPE rescaling factors tailored for the desired 2048k context length, followed by fine-tuning. Nevertheless, this strategy encounters significant hurdles due to the immense computational demands of training. In addition, our observations suggest that achieving effective fine-tuning with such high extension ratios is considerably difficult (refer to Appendix for details).

Fortunately, LongRoPE demonstrates efficacy when applied to both the base and fine-tuned variants of the extended LLM. Consequently, we propose a streamlined, incremental strategy that reaches the intended 2048k target utilizing only 1k fine-tuning steps, all conducted within a training sequence length of 256k.

$\Diamond$ \textit{LongRoPE search for expanding pre-trained LLMs to 256k}. Using LLaMA2 as a case study, we explore search strategies aimed at achieving context window sizes of 128k and 256k tokens. This corresponds to extension factors of 32$\times$ and 64$\times$ compared to the original setting.

$\Diamond$ \textit{Fine-tuning to 256k}. Next, the pre-trained LLM is fine-tuned to reach a 256k context window. The procedure involves an initial 400 steps of fine-tuning LLaMA2 using RoPE rescaled factors corresponding to 128k. Following this, the RoPE rescaled factors are swapped out for those appropriate for 256k on the obtained checkpoint, and an extra 600 steps of fine-tuning are performed. This staged approach demonstrates higher efficiency compared to immediately fine-tuning the model for 256k.

$\Diamond$ \textit{Scaling fine-tuned extended LLM to 2048k using LongRoPE search}. In the final stage, we conduct an additional search process on the LLM previously fine-tuned with a 256k context length. Through this approach, we achieve an exceptionally large context window of 2048k, all without requiring additional fine-tuning. This method yields a total extension factor of $512\times$.

\textbf{Recovery within reduced context windows}. When expanding the context window to a substantial size of 2048k, we observe diminished performance on tasks confined to the original context window length. This phenomenon is a recognized limitation associated with positional interpolation~\cite{pi}, which compels the position embeddings for the initial context range to be compressed into a smaller representational space in higher dimensions. As a result, the language model's efficacy is adversely impacted. Specifically, using an extension ratio of 512$\times$ leads to substantial congestion of the positions inside the original 4k context window.

To address this issue, we conduct an additional evolution search on the expanded LLM, specifically to fine-tune the RoPE rescale factors for shorter context lengths (such as 4k and 8k). For these shorter sequences, we limit the upper bound of the searched $\lambda$ values, since less positional interpolation is necessary. At inference time, the LLM adaptively modifies the RoPE rescale factors as needed.

\section{LongRoPE}

\label{sec:method}
Building upon these insights, we propose LongRoPE. This approach initially employs an effective search algorithm designed to leverage the two forms of non-uniformity, and subsequently applies this method to expand the context window of LLMs to more than 2 million tokens.

\subsection{Problem Formulation}

The presence of these two non-uniformities significantly expands the solution space and complicates the optimization process. To mitigate these challenges, we conceptualize the multidimensional non-uniform position interpolation optimization as a search problem.

Consider a LLM with a designated context window of size $L'$ and input documents $\mathbf{X}$, where each sequence $\mathbf{x} \in \mathbf{X}$ contains more tokens than $L'$. We represent the rotary angle in the RoPE embedding for the $i^{th}$ dimension at token position $n$ as $\frac{n}{\beta^i}$. The resulting optimization task can be stated as follows:

\begin{equation}
\begin{gathered}
		\label{eq:problem}

\underset {\mathbf{x}\in\mathbf{X};\, |\mathbf{x}|\ge L'}{ \operatorname {arg\,min} } \,  \mathcal{L} \left( \text{LLM}(\text{RoPE}, \mathbf{X}) \right), \quad \text{where}
\\
\scalebox{0.93}{
$\underset {\begin{subarray}{l} i=0,\cdots,\frac{d}{2}-1;\\ n\in[0,|\mathbf{x}|); \end{subarray}}{\text{RoPE}_{(n)}} = \left[ \cdots, \cos\left( \mathbb{I}(\hat{\lambda}_i, \hat{n}) \cdot \frac{n}{\beta^i} \right), \, \sin\left( \mathbb{I}(\hat{\lambda}_i, \hat{n}) \cdot \frac{n}{\beta^i} \right), \cdots \right] $}\\
\text{with} \quad \mathbb{I}(\hat{\lambda}_i,\hat{n}) = \begin{cases}
	1 & \text{if } n < \hat{n} \\
	\frac{1}{\lambda_i} & \text{if } n \geq \hat{n}
\end{cases}
\end{gathered}
\end{equation}

In this approach, we define a group of rescale factors, $\mathbb{I}(\hat{\lambda}_{i},\hat{n})$, to account for the two types of non-uniformities considered. Here, $\hat{\lambda_i}$ corresponds to non-uniformity along RoPE dimensions, while $\hat{n}$ refers to non-uniformity across token positions. More precisely, the rescale function $\mathbb{I}(\hat{\lambda}_{i},\hat{n})$ adjusts the rotation angle specifically for the $i^{th}$ dimension of RoPE, where $\hat\lambda_{i}$ is the applied scaling coefficient and $\hat{n}$ determines the threshold token position. For the initial token positions where $n < \hat{n}$, the factor $\hat\lambda_{i}$ remains inactive, resulting in the use of the standard RoPE angle $\frac{n}{\beta^i}$. When the token position satisfies $n \geq \hat{n}$, the rescale factor is incorporated.

Assuming a desired context window length of $L'$, our goal is to determine the set of optimal rescale factors ($\mathbb{I}(\hat{\lambda}_{0},\hat{n})$, $\mathbb{I}(\hat{\lambda}_{1},\hat{n})$, ..., $\mathbb{I}(\hat{\lambda}_{i},\hat{n})$, ...) across each RoPE dimension from $1$ through $d$. The aim is for the given $\text{LLM}$, equipped with these rescaled $\text{RoPE}$ embeddings, to attain the lowest possible next-token prediction loss $\mathcal{L}$ (or equivalently, the lowest perplexity) when provided with input sequences $\mathbf{X}$ whose length equals $L'$.

\subsection{Searching the Non-uniform Position Interpolation}

To address the challenge outlined in Eq.~\ref{eq:problem}, we propose an efficient yet straightforward approach that identifies the best RoPE rescale factors, enabling comprehensive utilization of multidimensional positional embedding non-uniformities.

\textbf{Search space}. Our approach establishes an expansive search space that accommodates both types of non-uniformity, as detailed in Table~\ref{tbl:searchspace}. In particular, we permit independent tuning of the rescale factor for each dimension within the RoPE embeddings. For simplicity in defining the search space, we search over variables $\lambda_i$ and $\hat{n}$ instead of directly optimizing $\mathbb{I}(\hat{\lambda}_{i},\hat{n})$, where we denote $\hat{\lambda}_{i}=1/\lambda_i$. As described in Table~\ref{tbl:searchspace}, the search range for $\lambda_i$ spans from a lower bound of 1.0 (representing standard extrapolation) up to an upper bound of $s\times1.25$ (allowing for broader interpolation than PI), with incremental steps of 0.01. Here, $s$ denotes the ratio for extending the target context window.

The parameter $\hat{n}$ determines how many of the earliest token positions preserve the unaltered position encoding, specifically by applying the original RoPE embedding rather than any interpolation. In practical terms, $\hat{n}$ is selected from the following set: \{0, 1, 2, 4, 8, 12, 16, 20, 24, 28, 32, 64, 128, 256\}. If $\hat{n}$ is set to 0, this implies that none of the token positions utilize the original embedding, and thus all positions are subjected to the optimized rescale factors found during the search process.

\textbf{Evolution-based search}. The range of possibilities outlined in Table~\ref{tbl:searchspace} encapsulates a vast number of positional interpolation configurations, making effective exploration highly challenging. For instance, choosing a $s=4\times$ extension results in $400^{128/2}\times14=4\times10^{167}$ potential options. As the extension ratio increases, the search space grows exponentially larger. To efficiently navigate this expansive landscape, we employ evolution search as described in~\cite{guo2020single} and incorporate two optimization strategies designed to significantly improve the efficiency of the search process. The complete workflow is depicted in Algorithm~\ref{alg:search}.

\textit{Optimized initial population generation}. Rather than selecting all $P$ rescale factors at random to form the initial population, we deliberately include the three RoPE rescale factors associated with PI, NTK, and YaRN as specific individuals at the outset. The remaining $P$-3 members of the population are created by introducing random mutations to these three base rescale factors, each with a probability $p$.

\textit{Monotonically non-decreasing constraint}. Once the initial population is produced, we evaluate the LLM perplexity for every candidate. This involves applying the selected RoPE rescale factors to the designated LLM and then determining the perplexity over the input $\mathbf{X}$. The best-performing $k$ candidates are chosen as parents for the next stage of evolution. Nonetheless, due to the enormous size of the search space, uninformed mutation and crossover can frequently result in unproductive explorations, thereby triggering superfluous perplexity evaluations. This inefficiency is especially pronounced for large $L'$, since each perplexity computation demands significant inference time.

To this end, we enforce a non-decreasing constraint on the sampled RoPE rescaling factors, such that $\lambda_i \le \lambda_{i+1}$ holds for all indices. We apply only those RoPE configurations that conform to this constraint during LLM perplexity evaluation, resulting in a substantial reduction in the search space. More precisely, we mandate that the sequence $\lambda_i$ should be monotonically increasing with respect to the RoPE dimension ($i = 0, \ldots, 63$). This requirement for dimensional monotonicity draws on insights from NTK theory~\cite{jacot2018neural,tancik2020fourier,ntk}, which indicates that lower-indexed (higher frequency) dimensions benefit from less interpolation (i.e., smaller $\lambda_i$ values), whereas higher-indexed (lower frequency) dimensions can accommodate greater interpolation (i.e., larger $\lambda_i$ values).

\begin{algorithm}[t]
	\small
	\caption{The search algorithm}
	\textbf{Input:} target LLM, input samples $\mathbf{X}$, population size $P$, mutation size $N_1$, crossover size $N_2$, max iterations $\mathcal{T}$, mutate probability $p$\\

	\begin{algorithmic}[1]
		\label{alg:search}
		\STATE Top-k $\leftarrow \phi$;	
		\STATE $\text{P}_0$ $\leftarrow$ \textit{Initialize\_population\_with\_optimization} ($P$, $\mathbf{X}$, $p$); 
		\FOR{$i$ from $1$ to $\mathcal{T}$}
		\STATE \textit{Compute\_perplexity} (LLM, $\text{P}_{i-1}$, $\mathbf{X}$);
		\STATE Top-k $\leftarrow$ \textit{Update\_Topk} (Top-k, $\text{P}_{i-1}$);
		\STATE $\text{P}_{mutation}$ $\leftarrow$ \textit{Mutation\_with\_mono\_constraint} (Top-k, $N_1$, $p$);
		\STATE $\text{P}_{crossover}$ $\leftarrow$ \textit{Crossover\_with\_mono\_constraint} (Top-k, $N_2$);
		\STATE $\text{P}_i$ $\leftarrow$ $\text{P}_{mutation}$ $\cup$ $\text{P}_{crossover}$ $\cup$ Top-k;
		\ENDFOR
		\STATE Return the member from Top-k exhibiting the lowest perplexity;
	\end{algorithmic}
\end{algorithm}

\textbf{8$\times$ extension without fine-tuning}. Leveraging our evolutionary search approach, we successfully determine non-uniform RoPE rescale factors that maintain crucial positions and dimensions, thereby reducing information loss caused by interpolation. As shown in Fig.\ref{fig:non-ft}, our strategy enables the expansion of LLaMA2’s context window from 4k to 32k tokens without requiring any fine-tuning. By comparison, alternative techniques like PI and non-uniform NTK and YaRN result in sharp increases in perplexity once the window is doubled.

\subsection{Extending LLM Context Window to 2048K}
\label{sec:extendto2048k}

\textbf{Progressive extension to 2048k}. Here, we present our approach for expanding the context window of pre-trained LLMs from the standard 4k up to beyond 2048k. As shown previously, utilizing our non-uniform positional interpolation strategy allows for an 8$\times$ increase in context length without necessitating fine-tuning. For substantially greater scaling (such as up to 512$\times$), fine-tuning becomes indispensable. One possible approach involves determining appropriate RoPE rescaling factors at the target 2048k context length followed by fine-tuning. Nevertheless, this strategy is often constrained by the immense computational resources demanded. Furthermore, our observations indicate that effectively fine-tuning LLMs with such large context expansions poses significant difficulties (refer to Appendix).

Fortunately, LongRoPE demonstrates efficacy when applied to both the base and fine-tuned variants of the extended LLM. To this end, we present a streamlined and incremental approach capable of reaching the target length of 2048k using only 1k fine-tuning iterations and a maximum training sequence length of 256k.

$\Diamond$ \textit{LongRoPE search for expanding pre-trained LLMs to 256k}. Using LLaMA2 as a case study, we perform a search aimed at increasing the context window to 128k and 256k. At these respective window sizes, the model undergoes an extension by factors of 32$\times$ and 64$\times$.

$\Diamond$ \textit{Fine-tuning to 256k}. In this phase, we adapt the pre-trained LLM to accommodate a context window size of 256k tokens. Initially, we conduct 400 fine-tuning steps on LLaMA2 utilizing RoPE rescaling parameters set for 128k. Following this, we switch the RoPE rescaling factors to target 256k using the previously obtained checkpoint, and continue with 600 further fine-tuning steps. This staged approach demonstrates greater efficiency compared to commencing fine-tuning directly at 256k.

$\Diamond$ \textit{Adapting fine-tuned extended LLM to 2048k via LongRoPE exploration}. In the final step, we conduct an additional exploration on the previously fine-tuned LLM, which had a 256k sequence length. As a result, we achieve a vast context window of 2048k, accomplishing this expansion without any further fine-tuning. The resulting extension rate reaches $512\times$.

\textbf{Recovery within the original context window}. Expanding the context window to an extensive length of 2048k results in degraded model performance when evaluated on the initial context window. This phenomenon is consistent with known limitations of positional interpolation~\cite{pi}, which compresses the position embeddings for the original context into a constrained segment of the embedding space, thereby impairing the language model’s efficacy. Notably, with an extension ratio of 512$\times$, the embeddings corresponding to positions in the original 4k context window become highly concentrated.

To address this, we conduct an additional evolutionary search on the expanded LLM to fine-tune the RoPE rescale factors specifically for shorter context lengths (such as 4k and 8k). The upper bound for the searched $\lambda$ is decreased because shorter context lengths necessitate less positional interpolation. In the inference phase, the LLM adaptively modifies the relevant RoPE rescale factors as needed.

 \section{Experiments}

\label{sec:exp}
\subsection{Setup}
\textbf{Evaluation Tasks and models}. LongRoPE is integrated with LLaMA2-7B and Mistral-7B. We assess their effectiveness in three main areas: (1) perplexity scores when handling lengthy input documents with extended-context LLMs; (2) performance on the Passkey retrieval task, designed to test a model’s proficiency in extracting a simple passkey embedded among extensive unrelated text; and (3) outcomes on conventional LLM evaluation benchmarks restricted to a 4096-token context window.

\textbf{Fine-tuning}. For the LLaMA2 model, fine-tuning is performed using a learning rate of 2e-5, which decays linearly, and a total batch size of 32 across all devices. Training proceeds for 400 iterations on the Redpajama~\cite{redpajama} dataset, partitioned into segments of 128k tokens, each delimited by BOS and EOS tokens. Subsequently, using the resulting checkpoint as a starting point, a further 600 steps are conducted to extend the context window to 256k tokens. Training the 128k context model utilizes 8 A100 GPUs with a distributed training architecture~\cite{cube}, whereas expanding to 256k tokens necessitates 16 A100 GPUs. For Mistral, training is performed with a fixed learning rate of 1e-6 and a global batch size of 64. In line with the approach in YaRN~\cite{yarn}, both the 128k and 256k models undergo 400 training steps on Together Computer’s Long-Data Collections~\cite{mistraldata} dataset with a 16k sequence length, making use of 4 A100 GPUs.

\textbf{Search}. For target window sizes up to 256k, the following configuration is utilized: $P$=64, $N_1$=$N_2$=16, $p$=0.3, and $\mathcal{T}$=40, with the top-32 individuals chosen for mutation or crossover at each generation. Perplexity evaluation is based on 5 randomly selected samples from the PG19 validation set, each sample meeting at least the target context length. For window sizes above 512k, population size along with mutation and crossover numbers are reduced by half. In this regime, perplexity is assessed using 3 randomly chosen instances from the Pile-Books3~\cite{pile} validation set.

\textbf{Baselines}. In order to achieve a 2048k context window, we fine-tuned our models using context lengths of 128k and 256k, resulting in the variants LongRoPE-2048k (ft=128k) and LongRoPE-2048k (ft=256k) for LLaMA2 and Mistral, respectively. Our evaluation includes these four models alongside leading context window extension baselines. Specifically, we consider open-source LLMs that have been further fine-tuned following positional interpolation methods such as PI, NTK, and YaRN. Among these are Together-32k~\cite{together}, Code LLaMA~\cite{codellama}, LongLoRA-full-FT-100k~\cite{longlora}, as well as YaRN-LLaMA and YaRN-Mistral~\cite{yarn}.

\subsection{Main Results}

\label{sec:mainresult}
\textbf{Language modeling on sequences up to 256k tokens}. Our initial analysis focuses on a comparison with leading extended LLM models when evaluated on sequences as long as 256k tokens. To assess model generalization, we employ the Proof-pile~\cite{pg19} and PG19~\cite{pile} test sets. Perplexity is measured across different context lengths, applying a sliding window of size 256. For the PG19 dataset, evaluations are conducted across the entire test split comprising 100 documents. For Proof-pile, we adopt the YaRN~\cite{yarn} protocol, sampling 10 instances at random, each containing a minimum of 128k tokens.

Table~\ref{tbl:proof-pile} and Table~\ref{tbl:pg19} present a comparison of perplexity scores for LLaMA2 and Mistral models enhanced using various interpolation techniques on the Proof-pile and PG19 datasets, respectively. Two main findings emerge from these results: \textbf{(1)} the extended models demonstrate a consistent decline in perplexity as evaluation lengths increase from 4k to 256k, indicating their improved ability to utilize extended contexts. \textbf{(2)} Remarkably, even when the context window is expanded by a factor of 16—a scenario where maintaining short-range performance is often problematic—our LongRoPE-2048k variants surpass the leading baselines for context lengths up to 256k.

\textbf{Long sequence language modeling beyond 2000k}. To assess performance on exceptionally lengthy documents, we utilize the Books3~\cite{pile} dataset. For efficient evaluation, we randomly sample 20 books, each longer than 2048k tokens, and employ a sliding window approach with a size of 256k.

As indicated in Table~\ref{tbl:books3}, LongRoPE is able to expand the context window of both LLaMA2-7B and Mistral-7B to 2048k tokens, while also maintaining perplexity levels that match or surpass those of baseline methods across shorter ranges of 8k-128k. A pronounced divergence in performance is observed between the 2048k versions of LLaMA2 and Mistral: although Mistral shows superior results over baselines at shorter sequence lengths, its perplexity surpasses 7 when evaluating beyond the 256k mark. The behavior of LLaMA2 conforms to expectations—its perplexity generally declines as context lengthens, aside from modest increases at the 1024k and 2048k intervals. Notably, for LLaMA2, fine-tuning LongRoPE-2048k with a window size of 256k yields better results than with 128k; this can be attributed to a smaller secondary extension ratio (namely, 8$\times$ rather than 16$\times$). In comparison, Mistral sees improved performance with a 128k fine-tuning window. This is chiefly due to our use of a 16k training length, following the YaRN configuration, for both the 128k and 256k fine-tuning of Mistral, which has implications for Mistral’s ability to further enlarge its context window post fine-tuning.

\textbf{Passkey retrieval}. Here, we investigate the model's performance in retrieving information from varying context window lengths during generation tasks. Specifically, we employ the synthetic passkey retrieval benchmark introduced by~\cite{passkey}. In this evaluation, the model must identify a randomly selected five-digit passkey embedded within a lengthy text passage. Details of the prompt template are provided in the appendix. We conduct ten runs of the retrieval task, each time inserting the passkey at a randomly selected position, uniformly distributed throughout the entire evaluation context.

Fig.~\ref{fig:passkey} presents a comparative analysis of retrieval accuracy with respect to baseline models. The performance of current models diminishes sharply to zero once the input exceeds 128k tokens. However, our LongRoPE-LLaMA2-2048k (ft=256k) model demonstrates robust retrieval capabilities, sustaining accuracy rates of at least 90\% across sequence lengths ranging from 4k to 2048k, despite the inherent complexity of extracting a passkey from token sequences numbering in the millions. Meanwhile, LongRoPE-Mistral-2048k (ft=128k) maintains perfect (100\%) retrieval accuracy up to 1800k tokens and then experiences a decrease to 60\% at 2048k. This result is consistent with the findings in Table~\ref{tbl:books3}, where there is a minor uptick in perplexity at the 2048k length.

\textbf{Standard benchmarks within original context window}. We assess the performance of LongRoPE-2048k models on their native context window by employing the Hugging Face Open LLM Leaderboard~\cite{huggingfaceboard} for both zero-shot and few-shot evaluations. Specifically, we conduct evaluations with 25-shot ARC-Challenge~\cite{allenai:arc}, 10-shot HellaSwag~\cite{zellers2019hellaswag}, 5-shot MMLU~\cite{mmlu}, and 0-shot TruthfulQA~\cite{lin2021truthfulqa}.

As presented in Table~\ref{tbl:commonsense}, our models deliver performance that is on par with previous results on the initial benchmark, which was tailored for a shorter context length, and even surpass the baseline Mistral by +0.5\% on TruthfulQA. Although LongRoPE-LLaMA2-2048k, which underwent fine-tuning at a context length of 256k, exhibits a modest decline in performance, it still stays within acceptable margins across the majority of tasks.

\subsection{Ablation Results}

\textbf{Effectiveness of the second positional interpolation}. Within our progressive extension framework, a second round of non-uniform positional interpolation is performed on fine-tuned, extended LLMs utilizing our search method. To assess its utility, we conduct experimental evaluations on our LLaMA2-256k model after fine-tuning. This model is then expanded to 512k, 1024k, and 2048k using both PI and YaRN techniques. According to the results reported in Table~\ref{tbl:secondsearch}, our approach with non-uniform positional interpolation maintains stable perplexity across these extensions. By comparison, applying PI and YaRN leads to a sharp rise in perplexity as the extension ratio increases.

\textbf{Recovery performance at short context windows.} In order to address the degradation in performance at reduced context lengths, we employ our search algorithm to re-optimize the RoPE factors for LongRoPE-2048k. In this process, we impose stricter upper bounds on the scale factors, limiting the extent of interpolation specifically for the 4k and 8k context ranges. Table~\ref{tbl:shortperformance} presents a perplexity assessment of LongRoPE-LLaMA2-2048k on Proof-pile for 4k and 8k contexts, in addition to average LLM benchmark accuracy. These results illustrate clear gains in model performance for shorter context lengths.

\textbf{Analysis on the two forms of non-uniformities}. To assess the impact of each non-uniformity on model performance, we perform ablation studies focusing on both types. Specifically, we conduct two experiments: (i) we extend LLaMA2-7B to short context lengths of 16k and 32k using various strategies—PI, optimizing over the RoPE dimension alone, and jointly optimizing both non-uniformities; (ii) we further scale our fine-tuned LLaMA2 model from a context window of 256k tokens to 2048k, employing the same experimental design. Perplexity is measured prior to any additional fine-tuning. The results summarized in Table~\ref{tbl:searchdimension} reveal that introducing non-uniformity into the RoPE dimension yields a pronounced perplexity reduction relative to the linear interpolation utilized in PI. While incorporating non-uniformity in token position enhances results at 16k and 32k, its benefit is not apparent at the 2048k sequence length, which may be attributed to the vast length involved. In such extreme cases, merely keeping the original tokens without applying interpolation seems ineffectual, and we propose to address this limitation in future investigations.

\section{Related Works}

In addition to methods based on position interpolation, this section discusses related works of other approaches. 

\textbf{Retrieval-based} methods incorporate an external memory component to store extensive historical context, employing retrieval mechanisms to access relevant documents during inference~\cite{longllama,wang2023augmenting,borgeaud2022improving}. Such frameworks often require deliberate and significant alterations to the LLM architectures. In contrast, our approach adopts a more streamlined strategy, involving only minimal changes to positional embeddings. Furthermore, our model extends its applicability to a broader range of long context tasks beyond retrieval, including long document summarization and few-shot learning.

\textbf{Attention-based context window extensions}. In addition to positional embedding interpolation, several studies have extended the input context length of the original LLM by altering attention mechanisms~\cite{han2023lminfinite, streamingllm,parallelwindow}. The central concept involves employing innovative attention masks to address the challenge of attention explosion arising from introducing new positions. These methods can function in tandem with positional interpolation strategies.

\textbf{Fine-tuning based approaches} investigate strategies for adapting pre-trained LLMs to extended context lengths through modified position embeddings. Prior efforts, including Code LLaMA~\cite{codellama}, LLaMA2 Long~\cite{xiong2023effective}, and ScaledRoPE~\cite{liu2023scaling}, employ significantly increased base values for RoPE and subsequently fine-tune models at the desired context window. In contrast, our approach accommodates a range of target context lengths and is capable of scaling to lengths surpassing 2M. More recent developments, such as LongLoRA~\cite{longlora} and PoSE~\cite{zhu2023pose}, address the considerable GPU requirements of long context fine-tuning (over 128k tokens) by introducing more efficient procedures. Our approach remains complementary to these resource-efficient fine-tuning methods.

\section{Conclusion}

In this study, we introduce LongRoPE, a technique that significantly increases the context length capacity of LLMs to an extraordinary 2048k, all while retaining their effectiveness within the original, more limited context window. Our approach harnesses two types of non-uniformities in RoPE positional embedding, identified via an efficient evolutionary search process. This methodology yields two main advantages: it offers advantageous initialization for subsequent fine-tuning and facilitates an 8$\times$ expansion of the context window even without fine-tuning. Furthermore, we put forth a progressive extension scheme, leveraging LLMs fine-tuned at a 256k context length to ultimately achieve a 2048k context window size, eliminating the need for additional fine-tuning. Comprehensive experimental results demonstrate the robustness of LongRoPE. We anticipate that the LongRoPE-2048k models will unlock a variety of novel long context use cases and encourage more research in this direction.

\section*{Broader Impacts} The objective of this paper is to contribute to the development of the Machine Learning field. While this research could have various implications for society, we do not identify any particular impacts that require special emphasis at this time.

	\balance

\end{document}